\section{Conclusion}
\label{sec:conclusion}

We presented 3DPS, a differentiable physically based point renderer for FMCW mmWave radar that synthesizes complex range profiles from novel sensor poses. The central design choice --- replacing optical-radiance primitives with oriented 3D points carrying ITU-R~P.2040 material parameters, evaluated by closed-form BSDFs and splatted through precomputed sensor PSFs --- eliminates the Monte Carlo sampling and post-FFT supervision that constrain prior radar-NVS methods, yielding a renderer that is simultaneously physics-faithful, complex-valued, and multi-viewpoint-tractable. On six outdoor ColoRadar scenes, 3DPS achieves \textbf{0.600} mean Pearson correlation on held-out novel-view $|\mathrm{RA}|$ --- $1.8\!\times$, $5.0\!\times$, and $7.4\!\times$ the strongest published baselines (RadarSplat, Radar Fields, DART) --- with physically interpretable per-point materials, in ${\sim}$3.6 minutes per scene over 8 training viewpoints (${\sim}$45$\times$ faster than mmIR's per-scene MC ray tracing cost of ${\sim}$160~min).

\textbf{Limitations and future work.} The current renderer evaluates single-bounce paths only, sufficient for the static outdoor scenes we target but not for indoor multipath-dominated environments where mmIR-style multi-bounce path tracing remains necessary. The geometric scaffold relies on a co-registered LiDAR cloud; relaxing this to LiDAR-free initialization is an open direction for purely radar-driven NVS. We supervise on $|\mathrm{RA}|$ magnitude for fair comparison with the three baselines, but a phase-aware loss exploiting the renderer's complex output should narrow the train--test gap further. Extension to dynamic scenes (range-Doppler-coherent rendering across moving objects), to downstream perception tasks (radar-trained detectors using NVS-augmented data), and to broader radar geometries beyond cascaded MIMO are natural next steps.
